\section{Corrections and reconciliation of the existing material}
\label{au:section}

The audit distinguishes a definite transcription error from research-status
updates and from questions that remain unresolved. The supplied patch makes
only the definite correction. The integration note proposes the status
changes separately.

\subsection{A reciprocal-Gamma coefficient contains a transcription error}

In the canonical chapter
\src{07-digamma-power-moments.tex}, equation
\src{dgp:eq:a5}, the final term of the displayed numerator is
\verb|+24\zeta\alpha_5|.
It should be \(24\zeta_5\), where the chapter uses
\(\zeta_j=\zeta(j)\). The corrected coefficient is
\begin{equation}\label{au:alpha5}
 \boxed{\displaystyle
 \alpha_5=
 \frac{\gamma^5-10\gamma^3\zeta(2)+20\gamma^2\zeta(3)
       +15\gamma\zeta(2)^2-30\gamma\zeta(4)
       -20\zeta(2)\zeta(3)+24\zeta(5)}{120}.}
\end{equation}
Indeed, the logarithm of the reciprocal Gamma function is
\[
 \log\frac1{\Gamma(1+u)}
   =\gamma u-\frac{\zeta(2)}2u^2+\frac{\zeta(3)}3u^3
          -\frac{\zeta(4)}4u^4+\frac{\zeta(5)}5u^5+O(u^6).
\]
Exponentiating and extracting \(u^5\) gives
\eqref{au:alpha5}. All other terms in the source's numerator agree with
this calculation. This is a transcription correction, not a new
reciprocal-Gamma identity. The companion patch was checked against the
pinned source with \src{git apply --check}; it was not applied to the
repository.

\subsection{The diagonal negative-order target was already completed}

Research question R7 in \emph{Tail Kernels and Boundary Products}
\cite{ProveItTailKernels} asks for the continued alternating-square law
at nonpositive orders, every order derivative, and the distinction between
ordinary summation and cutoff continuation. The incoming
\emph{Bilinear Tails and Determinants} report already supplies precisely
that diagonal result in
\src{sections/04_cutoff_jets.tex}, including its finite endpoint,
even-power expansion, contacts, translated cutoffs, and convergent
subtractions \cite{ProveItBilinear}.

The appropriate update is to cross-reference that proof and retain the
genuinely broader independent-order question. The latter was explicitly
left open in the same report and is completed by
Theorem~\ref{ct:contacts} here. The diagonal result is also an exact
computational compatibility check on the present bivariate coefficient
algorithm.

\subsection{The cubic Stieltjes target has a preserved general formula}

R9 of \emph{Tail Kernels and Boundary Products} asks for an explicit
holomorphic generating germ for cubic Stieltjes products with independent
indices and positive total argument-derivative order. Its stated first
target is
\[
 \FP\int_0^1\gamma_1(x)^2\gamma_0'(x)\dd x.
\]
The preserved \emph{Coincident Directional Zeta} continuation already
contains a general local completion for any number of factors and a finite
cubic Tornheim--zeta coordinate formula
\cite{ProveItCubic}. Its relevant source labels are
\src{cub:general-completion},
\src{cub:all-moments}, and \src{cub:global}.
Thus the existence of an explicit finite generating formula should be
marked as source-proved, subject to the project's normal integration and
review process. The further question of reducing selected retained
Tornheim order derivatives to a smaller arithmetic alphabet remains
meaningful. The supplied source does not prove general depth-one closure
or arithmetic independence.

The local mechanism was checked independently as follows. Write
\[
 H_{u,r}(x)=\partial_x^r\left(\zeta(1+u,x)-\frac1u\right).
\]
Its singular term is
\[
 (-1)^r(1+u)_r x^{-1-u-r},
\]
and its remainder is jointly holomorphic near \(u=x=0\). In a product
of three factors, choosing the singular term in a nonempty proper subset
\(S\) produces a resonance exactly at Taylor degree
\begin{equation}\label{au:resonant-degree}
  \sum_{i\in S}(1+r_i)-1
\end{equation}
of the remaining regular factors. The corresponding monomial is
\(A_S(\boldsymbol u)x^{-1-\sum_{i\in S}u_i}\). Its analytically
continued integral is
\(-A_S(\boldsymbol u)/\sum_{i\in S}u_i\), while all its order derivatives
at the resonance have fixed-coordinate finite part zero.

For \((r_1,r_2,r_3)=(0,0,1)\), the three singleton Taylor degrees are
\(0,0,1\); the three pair degrees are \(1,2,2\). These are all six
proper nonempty subsets. The full triple has no remaining regular factor
and has baseline exponent \(-4\), so it supplies no resonance through the
parameter origin. This checks the number and degrees of the local terms
required for the stated first target.

The correction uses the \emph{full analytic numerator}
\(A_S(\boldsymbol u)\) in each quotient. Replacing it by its restriction
to \(\sum_{i\in S}u_i=0\) removes the pole but changes transverse finite
coefficients. This is the same structural reason that the ambient contact
polynomials in Section~\ref{ct:section} must be differentiated before
restriction. The preserved cubic source retains these full numerators.
Its global Tornheim coordinates belong to the classical multiple-zeta
setting represented, for example, by Espinosa--Moll \cite{EspinosaMoll}.

\subsection{Results deliberately retained with their existing attribution}

The centered cubic trigamma family encountered during the audit is already
present in \emph{Quartic Spectral Reconstruction}
\cite{ProveItQuartic}, in the section containing
\src{qtm:C1}, \src{qtm:Cj}, \src{qtm:Rj}, and \src{qtm:G3}.
An independent finite summation reproduced the two exceptional
low-index corrections there. That family is not claimed as a contribution
of the present article.

Likewise, the degree bound for normalized spectral jets, the top local
polarization, and the all-difference formula already occur in
\emph{Arithmetic Spectral Jets}
\cite{ProveItArithmeticJets}, labels
\src{asj:thm:polynomial},
\src{asj:thm:top}, and
\src{asj:thm:all-differences}.
Equation~\eqref{sp:inherited-correction} is their specialized
normalization in the quadratic setting. The new ordinary density
representations should be integrated alongside this existing theorem,
not used to replace its attribution.

The unrestricted unequal-order reconstruction problem and the full
fixed-gap-six family remain open at the scope inspected here. None of the
contact or spectral identities in this article establishes their required
uniform nonvanishing. No further false theorem was established by this
targeted audit, and no claim is made that every proof in the complete
manuscript has been rechecked.
